\chapter{Experimentos y Resultados}

Este capitulo resume la corrida consolidada de la fase B2 y presenta resultados cuantitativos preliminares con foco en trazabilidad para paper y tesis.

\section{Configuracion experimental}

La corrida B2 se ejecuto en lotes y luego se consolidaron resultados en artefactos finales. Los archivos de referencia principales son:
\begin{itemize}
\item \texttt{results/opt\_benchmark\_b2\_full\_scale\_summary.csv}
\item \texttt{results/b2\_publication\_ranking\_final.csv}
\item \texttt{results/b2\_publication\_topk\_by\_family\_scenario.csv}
\end{itemize}

Las metricas reportadas fueron MAE, RMSE, DTW y SNR. Para cada combinacion de dataset-escenario-metodo se consolidaron estadisticos de tendencia central y variabilidad.

\section{Resultados}

En los escenarios analizados no aparece un ganador unico universal. La evidencia B3/B4
resume el rendimiento por familia (baseline, GSP y TV/tiempo), mostrando un compromiso
entre error puntual (MAE) y alineamiento temporal (DTW).

\begin{table}[h]
\centering
\caption{Resumen final REP-01 (B1+B2 consolidados).}
\begin{tabular}{lcccc}
\hline
Grupo & MAE $\downarrow$ & DTW $\downarrow$ & RMSE $\downarrow$ & Metodo top \\
\hline
Baseline & 0.071270 & 0.521934 & 0.162150 & spherical\_spline \\
TV/Tiempo & 0.072925 & 0.491819 & 0.167430 & tv \\
GSP & 0.073695 & 0.498893 & 0.168976 & gsmooth \\
\hline
\end{tabular}
\label{tab:tesis_rep01_final}
\end{table}

Ademas, se cerro STAT-02 con pruebas de significancia sobre la corrida full-scale:
el contraste por familias instantaneo vs TV/tiempo rechazo la hipotesis nula tanto en
MAE como en DTW (Wilcoxon pareado, $\alpha=0.05$). A nivel de metodos, TRSS vs BGSRP
mostro diferencia significativa en MAE, mientras que TRSS vs GSP no mostro diferencia
significativa en esa metrica.

Tambien se verifico consistencia operativa en la corrida B2, ya que el registro consolidado de warnings no reporta entradas en \texttt{results/opt\_benchmark\_b2\_full\_scale\_warnings\_registry.csv}. Esto no prueba perfeccion del pipeline, pero si indica que la corrida larga fue estable en terminos practicos.

Finalmente, para BGSRP se mantiene una nota de cautela: el archivo \texttt{results/b2\_publication\_bgsrp\_gap\_residual.csv} reporta un gap residual en comparacion proxy, aceptado en B2 pero todavia pendiente de cierre estricto 1:1 con MATLAB/GSPBox.
